\chapter{Introduction}
\label{ch:introduction}

\section{A flight is a transport process with an energy budget}

A soaring pilot travels by repeatedly exchanging height for distance and recovering height
in rising air. The atmosphere supplies the energy, but the pilot must locate that supply
and decide how to use it. A long displacement is therefore the outcome of several coupled
processes: motion through the air, encounters with lift, decisions made with incomplete
information, and constraints imposed by wind and terrain. Because the recorded trajectory
contains all these contributions, inferring a stochastic description requires more than
fitting its geometry: the description must also preserve the physical mechanisms that
allow the pilot to remain airborne.

Two elementary relations explain why neither an unconstrained random walk nor a purely
geometrical description is sufficient. Writing $\mathbf r$ for horizontal position and $z$
for altitude, a local kinematic description is
\begin{align}
  \dot{\mathbf r}(t)
    &=\mathbf v_{\mathrm{air}}(t)+\mathbf w_h(\mathbf r(t),z(t),t),
    \label{eq:intro-groundspeed}\\
  \dot z(t)&\simeq w_z(\mathbf r(t),z(t),t)
                  -s\bigl(|\mathbf v_{\mathrm{air}}(t)|,\phi(t)\bigr).
    \label{eq:intro-height}
\end{align}
Here $\mathbf w_h$ and $w_z$ are the horizontal and vertical air velocities, $s>0$ is the
sink rate relative to the air, and $\phi$ is bank angle. The second relation is a
quasi-steady approximation: it makes explicit the competition between lift and aerodynamic
sink without attempting to model every manoeuvre. Altitude is therefore a finite resource,
and a transport model must account for its use and replenishment if its trajectories are
to represent sustained flight.

This physical picture suggests three descriptions of behaviour: directed
\emph{transition}, exploratory \emph{search}, and \emph{climb} in lift. They are hypotheses
about the organisation of motion, rather than labels directly supplied by a recorder.
A pilot can find lift during a transition, climb without completing a circle, or search
while losing height. Determining which distinctions the data support is part of the
scientific problem.

Soaring control has been studied through the problem of locating and exploiting thermals,
including adaptive strategies based on local information~\cite{reddy2016}. The present
work asks a complementary question: what statistical laws emerge over entire flights, and
which mechanisms can account for those laws? A direct empirical precedent is the analysis
of soaring transport and flight phases by Vilpellet et al.~\cite{vilpellet2026}. This
thesis revisits that connection using the FFVL archive, explicit data-quality controls,
and a separation between measured properties and assumptions used to interpret them.

\section{The questions and the evidence needed to answer them}

\subsection{Which stochastic models survive the observations?}

A first observable is the squared displacement over an interval $\tau$,
\begin{equation}
  M_2(\tau)=\left\langle\left|\mathbf r(t+\tau)-\mathbf r(t)\right|^2\right\rangle.
  \label{eq:intro-msd}
\end{equation}
Ordinary diffusion without drift has $M_2\propto\tau$; directed motion can produce a ballistic law
$M_2\propto\tau^2$. An intermediate slope is informative, but it does not identify a
mechanism, since correlated Gaussian motion, persistent motion observed during a crossover,
and mixtures of flight behaviours can produce similar second moments. Random-walk
frameworks offer competing explanations of anomalous transport~\cite{metzler2000}.
The purpose here is to discriminate between such explanations using more than one curve.

Chapter~\ref{ch:global} therefore measures displacement quantiles, the spectrum of
moments, the shape of the increment distribution, directional anisotropy and velocity
correlations alongside the MSD. Their roles are complementary: quantile ratios test
whether the distribution changes shape as the lag increases, while higher moments
compare the growth of ordinary and large displacements; velocity correlations resolve
persistence in time. Open and closed tasks help expose the effect of the flight's intended
geometry. A model is challenged where its predictions disagree with these observations
under a comparable measurement protocol.

This logic also sets the limit of an exclusion. A discrepancy with a homogeneous Gaussian
model for pooled flights need not exclude Gaussian fluctuations within a flight: different
conditions can create a non-Gaussian mixture. Similarly, evidence against a particular
constant-speed renewal L\'evy walk is not a theorem against every process with finite-speed
legs~\cite{zaburdaev2015}. The chapter specifies the hypotheses being tested and leaves
unsupported exclusions open.

\subsection{How do the flight phases generate global transport?}

Although an unsegmented flight supplies constraints on the whole process, it does not
identify which behaviour generates a long displacement, a change of heading or a period
of persistence.
Chapter~\ref{ch:phases} addresses this missing link with a probabilistic segmentation into
transition, search and climb. It states the mathematical model, displays its behaviour on
readable trajectory excerpts, and distinguishes internal consistency from accuracy against
independent annotations.

The next step will be a chapter devoted to observables conditioned on phase: episode
durations, distance and height changes, speed and turning distributions, displacement
scaling, and dependencies between successive episodes. A useful model must reproduce both
these conditional statistics and the global ones, a stronger requirement than generating
a visually plausible sequence of coloured flight segments. Windows that cross
phase boundaries must also be retained when reconstructing global observables.

\subsection{Does information from other pilots change the motion?}

A nearby pilot can reveal lift before a follower reaches it. This creates a possible
information channel that is absent from a model of independent searchers. A later chapter
will compare solitary and group flight using both global and phase-conditioned observables.
The operational definition of a group will concern simultaneous proximity and plausible
visibility, not membership of a competition. A competition can contain isolated pilots,
and pilots outside a competition can travel together.

The comparison will ask where differences occur: in time spent searching, climb selection,
transition direction, episode persistence, or the efficiency of converting altitude into
distance. Site, weather, wing and pilot characteristics must be controlled because they can
produce the same differences without social interaction. Incomplete recording of nearby
pilots will limit what can be called solitary flight. These are conditions for an
interpretable comparison, not effects assumed in advance.

\section{From empirical constraints to a physical model}
\label{sec:intro-method}

The thesis combines empirical analysis, mathematical modelling and simulation, with a
distinct role for each. Data analysis establishes the constraints that a candidate process
must satisfy; mathematics then makes that process precise and derives its predictions,
which Monte Carlo simulation tests under finite flight durations, heterogeneous environments
and the actual observation protocol.

Model construction will follow the global analysis and the analysis conditioned on phase.
Possible ingredients include finite-speed transitions, lift encounters, distributions of
episode duration, and dependencies between successive directions or episodes. Renewal
independence is one hypothesis to test; correlated random walks provide alternatives when
that independence fails~\cite{tejedor2010,schulz2013}. The choice among these ingredients must be justified by the observations and an explicit
physical mechanism, rather than fixed by the outline of the thesis. Parameters fitted on one set of flights will be checked on
separate flights or days, using several observables jointly.

A reference model without wind and terrain may help isolate mechanisms, but obtaining it
requires an inference about the environment: mean ground velocity alone does not determine
wind, and rotating a trajectory into its principal axes does not remove orography. Artificially
making a covariance matrix isotropic changes the measured object; it does not recreate a
flight in still air. The analysis consequently preserves the motion in geographical
coordinates and studies environmental strata before asking which properties could survive
in a common reference model.

A further part of the programme concerns the spatial organisation of thermals and route
choice. It will first estimate the two-dimensional distribution of \emph{visited} lift
regions, accounting for where pilots have searched. Recovering the distribution of
available thermals requires an observation model: unvisited terrain is not evidence of
absent lift. A mathematical description of thermal availability and strength could then
support separate optimisation problems, such as minimum-time travel between two locations,
a closed route, or maximum distance within a specified flight budget. Wind and, subsequently,
possible geometrical constraints from terrain will be introduced in stages. A spatial
density alone cannot determine an optimal path without assumptions about thermal strength,
spatial dependence, altitude and the pilot's information.

\section{Structure and current scope}

Figure~\ref{fig:thesis-structure} shows how the measurements support the later modelling.
The present document develops the data, global transport and segmentation chapters. The
chapters on conditional observables, social flight, stochastic modelling and route
optimisation form the research programme; their results are not claimed here. The immediate
priority is to consolidate Chapters~\ref{ch:global} and~\ref{ch:phases}, then undertake the
solo--group comparison with the necessary observables conditioned on phase.

\begin{figure}[H]
  \centering
  \begin{tikzpicture}[schematic, node distance=5.5mm,
      % Present chapters read as solid, filled blue; the research programme as
      % dashed grey, so the two groups separate before any word is read.
      box/.style={rounded corners=2pt, align=center, text width=12.6cm,
                  inner xsep=8pt, inner ysep=7pt, line width=0.7pt},
      done/.style={box, draw=figblue, fill=figblue!6},
      future/.style={box, draw=black!45, dashed, fill=black!3},
    ]
    % A centred box breaks words when a line is a touch too long; the flow chart
    % is short enough to forbid that outright.
    \hyphenpenalty=10000 \exhyphenpenalty=10000
    \node[done] (data) {{\bfseries Chapter 2: trustworthy trajectories}\\[2pt]
      {\footnotesize\color{black!70} Recorder channels, cleaning, coverage and retained uncertainty}};
    \node[done, below=of data] (global) {{\bfseries Chapter 3: global observables}\\[2pt]
      {\footnotesize\color{black!70} Scaling, distribution shape, memory and environmental geometry}};
    \node[done, below=of global] (phases) {{\bfseries Chapter 4: flight phases}\\[2pt]
      {\footnotesize\color{black!70} Explicit segmentation model, diagnostic figures and independent validation}};
    \node[future, below=of phases] (conditional)
      {{\bfseries Observables conditioned on phase; solo--group comparison}\\[2pt]
      {\footnotesize\color{black!70} Mechanisms within episodes, dependencies between episodes, social information}};
    \node[future, below=of conditional] (model)
      {{\bfseries Stochastic model and Monte Carlo}\\[2pt]
      {\footnotesize\color{black!70} Joint reproduction of global and conditional observables on held-out data}};
    \node[future, below=of model] (routes)
      {{\bfseries Thermal landscapes and route optimisation}\\[2pt]
      {\footnotesize\color{black!70} Lift availability, wind, altitude budget and possible terrain constraints}};
    \foreach \a/\b in {data/global, global/phases, phases/conditional,
                       conditional/model, model/routes}
      \draw[->, black!45, line width=0.7pt] (\a) -- (\b);
    % Group labels in the left margin of the float, set against the brace they span.
    \draw[decorate, decoration={brace, amplitude=4pt}, black!45, line width=0.6pt]
      ($(data.north west)+(-3mm,0)$) -- ($(phases.south west)+(-3mm,0)$);
    \node[fig note, rotate=90, anchor=south]
      at ($($(data.north west)!0.5!(phases.south west)$)+(-6mm,0)$) {this document};
    \draw[decorate, decoration={brace, amplitude=4pt}, black!45, line width=0.6pt]
      ($(conditional.north west)+(-3mm,0)$) -- ($(routes.south west)+(-3mm,0)$);
    \node[fig note, rotate=90, anchor=south]
      at ($($(conditional.north west)!0.5!(routes.south west)$)+(-6mm,0)$) {subsequent work};
  \end{tikzpicture}
  \caption{Logical structure of the thesis. Solid boxes are developed in this document;
  dashed boxes describe subsequent work. Model comparison must return to the measured
  observables, not stop at a fitted trajectory or a single exponent.}
  \label{fig:thesis-structure}
\end{figure}

Chapter~\ref{ch:dataset} follows the records from acquisition to retained trajectories,
including the treatment of missing or inconsistent channels. Chapter~\ref{ch:global}
then asks what can be learned without assigning phases. Chapter~\ref{ch:phases} introduces
that assignment and identifies what must be checked before its labels can support physical
claims. The two appendices state the spectral conventions and the geodetic
transformation used by the measurements.

\section{Chapter summary}
The thesis asks which stochastic description can account for soaring trajectories while
respecting the physical roles of lift, travel and environmental constraints. The route
from observations to a model begins with reliable trajectories and global transport
statistics, then uses validated flight phases to identify the mechanisms mixed by those
statistics. Comparisons between solitary and interacting flights, followed by simulation
and thermal-landscape modelling, extend this programme. Their conclusions remain future
work: the present chapters establish the empirical measurements and validation needed
to make those comparisons meaningful.
